\chapter{Category Theory in Brief}
\label{app:cat}

This appendix collects the definitions used in the main text. Standard references are~\cite{maclane1998,awodey2010}.

\section*{Categories and functors}

A \emph{category} $\mathbf{C}$ consists of a class of objects, for each ordered pair of objects $x,y$ a set $\Hom(x,y)$ of morphisms, an associative composition $\circ:\Hom(y,z)\times\Hom(x,y)\to\Hom(x,z)$, and for each object an identity morphism $\id_x$ satisfying $f\circ\id_x=f=\id_y\circ f$. A category with one object is the same as a monoid.

A \emph{functor} $F:\mathbf{C}\to\mathbf{D}$ assigns to each object $x$ an object $Fx$ and to each morphism $f:x\to y$ a morphism $Ff:Fx\to Fy$ with $F(g\circ f)=Fg\circ Ff$ and $F\id_x=\id_{Fx}$. It is \emph{faithful} if it is injective on each $\Hom(x,y)$ and \emph{full} if it is surjective on each. A contravariant functor reverses the direction of morphisms and is a functor $\mathbf{C}^{\op}\to\mathbf{D}$.

A \emph{natural transformation} $\eta:F\Rightarrow G$ between functors $F,G:\mathbf{C}\to\mathbf{D}$ assigns to each object $x$ a morphism $\eta_x:Fx\to Gx$ such that $\eta_y\circ Ff=Gf\circ\eta_x$ for every $f:x\to y$.

\section*{Limits and fiber products}

Given functors $F:\mathbf{A}\to\mathbf{C}$ and $G:\mathbf{B}\to\mathbf{C}$, the \emph{fiber product} $\mathbf{A}\times_{\mathbf{C}}\mathbf{B}$ has as objects pairs $(a,b)$ with $Fa=Gb$ and as morphisms pairs $(f,g)$ with $Ff=Gg$. It is the limit of the cospan $\mathbf{A}\to\mathbf{C}\leftarrow\mathbf{B}$ in the category of small categories.

\section*{Monoidal categories}

A \emph{monoidal category} is a category with a bifunctor $\otimes$ and a unit object $\mathbf{1}$ together with natural isomorphisms for associativity and unit, satisfying the pentagon and triangle coherence conditions. The \emph{interchange law} $(g\circ f)\otimes(g'\circ f')=(g\otimes g')\circ(f\otimes f')$ is part of the statement that $\otimes$ is a functor on the product category. A \emph{lax monoidal functor} $F$ between monoidal categories comes with a morphism $Fx\otimes Fy\to F(x\otimes y)$ natural in $x,y$ and a morphism $\mathbf{1}\to F\mathbf{1}$ satisfying coherence. When the morphisms are identities or isomorphisms the functor is \emph{strong}.

\section*{Adjunctions and Galois connections}

Functors $F:\mathbf{C}\to\mathbf{D}$ and $G:\mathbf{D}\to\mathbf{C}$ are \emph{adjoint}, written $F\dashv G$, if there is a bijection $\Hom_{\mathbf{D}}(Fx,y)\cong\Hom_{\mathbf{C}}(x,Gy)$ natural in $x$ and $y$. When $\mathbf{C}$ and $\mathbf{D}$ are posets, this says $Fx\le y\iff x\le Gy$, which is a monotone Galois connection. The antitone connection of Chapter~\ref{ch:scripts} is the same condition with one of the orders reversed.
